% Mathématiques 6e — cours V3
% Données, tableaux et graphiques

\section{Objectifs}

\begin{itemize}
\item Planifier une enquête simple.
\item Choisir les données à recueillir en fonction d'une question.
\item Recueillir des observations ou réaliser des mesures.
\item Consigner les données dans un tableau clair.
\item Construire un tableau simple pour présenter des données.
\item Lire un tableau, un diagramme en barres, un diagramme circulaire ou une courbe dans des situations immédiates.
\item Filtrer un tableau selon un critère pour effectuer un choix.
\item Lire une représentation graphique avec esprit critique.
\end{itemize}

\section{1. Une donnée répond à une question}

Une donnée n'est pas un nombre isolé sans contexte.

Elle est recueillie pour répondre à une question.

\begin{exemple}
On souhaite étudier les modes de déplacement utilisés par les élèves pour venir au collège.

La question peut être :
\[
\text{« Quel moyen de transport utilises-tu le plus souvent ? »}
\]

Les réponses possibles peuvent être : marche, vélo, bus, voiture ou autre.
\end{exemple}

Avant toute collecte, il faut donc déterminer précisément ce que l'on veut savoir.

\section{2. Planifier une enquête}

Une enquête peut être organisée en plusieurs étapes.

\coursefigure{/images/mathematiques/6eme/donnees-enquete-tableau.svg}{Schéma des étapes d'une enquête et exemple de tableau de données.}{Une enquête efficace relie la question, la collecte, l'organisation des données et l'interprétation finale.}

\begin{methode}
Pour préparer une enquête :
\begin{enumerate}
\item formuler une question précise ;
\item définir les personnes ou objets étudiés ;
\item choisir les informations à recueillir ;
\item décider comment noter les réponses ;
\item recueillir les données ;
\item organiser les résultats ;
\item interpréter sans dépasser ce que les données permettent d'affirmer.
\end{enumerate}
\end{methode}

Une question vague produit souvent des réponses difficiles à comparer.

\section{3. Population et individus étudiés}

La \textbf{population étudiée} est l'ensemble des personnes ou objets sur lesquels porte l'enquête.

Chaque élément de cette population est un individu de l'étude.

\begin{exemple}
Si l'on interroge les $28$ élèves d'une classe sur leur temps de trajet, la population de l'enquête est constituée de ces $28$ élèves.
\end{exemple}

Les conclusions concernent d'abord cette population.

Interroger une seule classe ne permet pas automatiquement de conclure pour tous les collégiens de France.

\section{4. Caractères observés}

Une enquête peut porter sur différents types d'informations.

\begin{itemize}
\item Une catégorie : moyen de transport, couleur préférée, activité choisie.
\item Une mesure : taille, durée, distance, température.
\item Un nombre : quantité de livres lus, nombre de trajets effectués.
\end{itemize}

L'information relevée doit être définie avant la collecte pour éviter des réponses ambiguës.

\section{5. Réaliser des mesures}

Certaines données ne sont pas déclarées mais mesurées.

On peut par exemple mesurer :
\[
\text{une longueur},\qquad
\text{une durée},\qquad
\text{une température}.
\]

Une mesure doit toujours être accompagnée de son unité.

\begin{attention}
Un tableau contenant des mesures sans unité est incomplet. Les valeurs $12$, $15$ et $18$ ne signifient pas la même chose selon qu'elles représentent des secondes, des centimètres ou des degrés Celsius.
\end{attention}

\section{6. Consigner les données}

Pendant la collecte, il faut noter les informations de manière régulière.

On peut utiliser un tableau avec une ligne par individu et une colonne par information.

\begin{exemple}
Pour une enquête sur le trajet vers le collège, on peut utiliser les colonnes :
\[
\text{Élève},\qquad
\text{moyen de transport},\qquad
\text{durée du trajet}.
\]

Cette organisation permet ensuite de filtrer ou de comparer les réponses.
\end{exemple}

\section{7. Construire un tableau simple}

Un tableau doit comporter des intitulés compréhensibles.

\begin{tabular}{c|c|c}
Élève & Transport & Durée en minutes \\
A & Vélo & $12$ \\
B & Bus & $25$ \\
C & Marche & $9$ \\
D & Vélo & $15$ \\
\end{tabular}

Les intitulés précisent ce que représentent les colonnes.

L'unité « minutes » n'a pas besoin d'être répétée dans chaque cellule lorsqu'elle apparaît clairement dans l'en-tête.

\section{8. Lire un tableau}

Lire un tableau consiste à retrouver une information à l'intersection d'une ligne et d'une colonne.

Dans le tableau précédent, la durée du trajet de l'élève C est :
\[
9\ \text{min}.
\]

Le moyen de transport de l'élève D est :
\[
\text{vélo}.
\]

\begin{methode}
Avant de lire une cellule :
\begin{enumerate}
\item repérer la bonne ligne ;
\item repérer la bonne colonne ;
\item vérifier l'intitulé et l'unité ;
\item seulement ensuite lire la valeur.
\end{enumerate}
\end{methode}

\section{9. Regrouper des réponses}

Des données individuelles peuvent être regroupées par catégorie.

Supposons que, sur $30$ élèves :
\[
12
\]
viennent en bus,
\[
8
\]
à vélo,
\[
6
\]
à pied et
\[
4
\]
en voiture.

On contrôle :
\[
12+8+6+4=30.
\]

Ce contrôle permet de vérifier qu'aucune réponse n'a été oubliée ou comptée deux fois.

\section{10. Effectif d'une catégorie}

L'effectif d'une catégorie est le nombre d'individus appartenant à cette catégorie.

Dans l'exemple précédent, l'effectif « vélo » vaut :
\[
8.
\]

Le mot effectif permet de résumer une collecte, mais le travail essentiel reste l'organisation correcte des données qui ont permis de l'obtenir.

\section{11. Diagramme en barres}

Un diagramme en barres représente des catégories par des barres dont la hauteur ou la longueur correspond à une valeur.

Pour le lire, on vérifie :
\begin{itemize}
\item le titre ;
\item les catégories ;
\item l'axe des valeurs ;
\item l'échelle ;
\item l'unité éventuelle.
\end{itemize}

Deux barres deux fois plus hautes ne signifient « deux fois plus » que si l'échelle utilisée le permet réellement.

\section{12. Diagramme circulaire}

Un diagramme circulaire représente la répartition d'un ensemble.

Le disque entier représente le total.

Une partie plus grande du disque correspond à une catégorie plus importante.

\begin{remarque}
En 6e, la lecture immédiate d'un diagramme circulaire peut être travaillée sans transformer systématiquement chaque secteur en calcul d'angle.
\end{remarque}

\section{13. Courbe d'évolution}

Une courbe permet souvent de représenter une grandeur qui évolue.

Par exemple, une température relevée toutes les heures peut être placée sur un graphique.

On peut alors répondre à des questions comme :
\begin{itemize}
\item à quel instant la température est-elle maximale ?
\item entre quels instants augmente-t-elle ?
\item quelle valeur lit-on à un instant donné ?
\end{itemize}

L'axe horizontal et l'axe vertical doivent être lus séparément.

\section{14. Filtrer un tableau}

Filtrer signifie ne conserver que les lignes qui respectent un critère.

Supposons que l'on cherche les élèves dont le trajet dure moins de :
\[
15\ \text{min}.
\]

On examine la colonne « durée » et on conserve seulement les lignes vérifiant :
\[
d<15.
\]

Le filtre transforme un grand tableau en une liste correspondant à une question précise.

\section{15. Filtrer avec plusieurs critères}

On peut imposer plusieurs conditions.

Par exemple :
\[
\text{transport = vélo}
\]
et
\[
d\le15\ \text{min}.
\]

Il faut alors conserver uniquement les lignes qui satisfont \textbf{les deux conditions}.

\begin{exemple}
Dans un tableau, trois élèves viennent à vélo. Leurs durées sont :
\[
12,\quad15,\quad19\ \text{min}.
\]

Avec le critère $d\le15\ \text{min}$, deux élèves sont retenus.
\end{exemple}

\section{16. Faire un choix à partir d'un filtre}

Le filtrage sert souvent à prendre une décision.

On peut chercher :
\begin{itemize}
\item les produits coûtant moins d'un certain prix ;
\item les trajets durant moins d'une certaine durée ;
\item les communes répondant à plusieurs critères ;
\item les mesures situées dans un intervalle fixé.
\end{itemize}

\begin{attention}
Le tableau ne décide pas à notre place. Il permet d'identifier les éléments qui respectent les critères choisis.
\end{attention}

\section{17. Représentation graphique et lecture critique}

Un graphique peut donner une impression visuelle très forte.

\coursefigure{/images/mathematiques/6eme/donnees-graphique-echelle.svg}{Deux diagrammes en barres utilisant des échelles verticales différentes.}{Des valeurs proches peuvent sembler très différentes si l'échelle est resserrée ; l'axe doit être lu avant toute interprétation.}

Pour interpréter correctement un graphique, il faut examiner :
\begin{itemize}
\item le point de départ de l'axe ;
\item la valeur d'un intervalle ;
\item la régularité de l'échelle ;
\item les unités ;
\item la période étudiée ;
\item la population concernée.
\end{itemize}

\section{18. Axe tronqué}

Un axe vertical n'est pas obligé de commencer à zéro.

Cela peut être utile pour voir de petites variations.

Mais cela peut aussi exagérer visuellement une différence.

\begin{exemple}
Deux valeurs valent :
\[
98
\]
et
\[
100.
\]

La différence est seulement :
\[
100-98=2.
\]

Sur un axe allant de $0$ à $120$, les barres seront proches.

Sur un axe allant de $97$ à $101$, elles sembleront beaucoup plus différentes.
\end{exemple}

Il faut donc lire les nombres avant de conclure à partir de l'apparence.

\section{19. Une enquête peut être biaisée}

Les résultats dépendent de la manière dont les données ont été recueillies.

\begin{exemple}
Pour connaître le sport préféré de tous les élèves d'un collège, interroger uniquement les membres de l'association sportive risque de produire une enquête peu représentative.
\end{exemple}

En 6e, l'objectif n'est pas de formaliser toute la théorie des sondages, mais de développer un premier réflexe critique :
\[
\text{Qui a été interrogé ?}
\]

\section{20. Exemple développé : organiser une enquête}

On veut connaître le temps de trajet des élèves d'une classe.

\begin{methode}
Une démarche possible est :
\begin{enumerate}
\item définir la population : les élèves de la classe ;
\item demander à chacun la durée habituelle de son trajet ;
\item imposer une unité commune, par exemple la minute ;
\item consigner une ligne par élève ;
\item vérifier que le nombre de lignes correspond au nombre de réponses ;
\item filtrer ensuite selon une question, par exemple $d\le10\ \text{min}$.
\end{enumerate}
\end{methode}

Si $28$ élèves répondent, le tableau final doit contenir :
\[
28
\]
lignes de données, sauf réponse manquante explicitement signalée.

\section{21. Exemple développé : choisir à partir d'un tableau}

Trois trajets sont décrits ainsi :

\begin{tabular}{c|c|c}
Trajet & Durée & Correspondance \\
A & $35\ \text{min}$ & non \\
B & $28\ \text{min}$ & oui \\
C & $32\ \text{min}$ & non \\
\end{tabular}

On cherche un trajet :
\begin{itemize}
\item sans correspondance ;
\item de durée inférieure à $34\ \text{min}$.
\end{itemize}

Le trajet A ne convient pas car :
\[
35>34.
\]

Le trajet B ne convient pas car il comporte une correspondance.

Le trajet C vérifie les deux critères.

Ainsi, le filtre retient :
\[
\boxed{\text{trajet C}}.
\]

\section{22. Méthode générale}

\begin{methode}
Pour travailler sur des données :
\begin{enumerate}
\item comprendre la question posée ;
\item identifier la population et les informations utiles ;
\item recueillir ou mesurer proprement ;
\item construire un tableau avec des intitulés et unités ;
\item contrôler les données ;
\item représenter si cela aide à répondre à la question ;
\item filtrer selon les critères demandés ;
\item interpréter en tenant compte de l'échelle et du contexte.
\end{enumerate}
\end{methode}

\section{23. Erreurs fréquentes}

\begin{itemize}
\item Commencer à construire un graphique avant d'avoir organisé les données.
\item Oublier les unités dans un tableau de mesures.
\item Lire une cellule sans vérifier sa ligne et sa colonne.
\item Confondre la hauteur visuelle d'une barre avec sa valeur sans lire l'échelle.
\item Filtrer sur un seul critère lorsqu'il en faut deux.
\item Tirer une conclusion sur une population beaucoup plus large que celle réellement étudiée.
\item Croire qu'un graphique est neutre sans examiner ses axes.
\end{itemize}

\section{À retenir}

\begin{aretenir}
En 6e, le travail sur les données commence par une question et une collecte organisée.

Il faut savoir planifier une enquête, recueillir des données ou effectuer des mesures, puis les consigner dans un tableau clair.

Un tableau permet ensuite de lire, comparer et \textbf{filtrer} les données selon un critère.

Les diagrammes et courbes sont des représentations : leur titre, leurs axes, leurs unités et leur échelle doivent être vérifiés avant toute interprétation.

Une représentation peut faciliter la compréhension, mais elle ne remplace jamais la lecture des données elles-mêmes.
\end{aretenir}
